\documentclass{situationist-preprint}

\title{The Flicker of Freedom Objection}
\msid{ACT-02-b}
\msfamily{ACT-02}
\msdiscipline{Action, Agency and Free Will}
\mstype{Companion manuscript}
\mskeywords{Frankfurt cases; flicker of freedom; robustness; alternative possibilities; moral responsibility}

\begin{document}
\maketitle

\begin{abstract}
The flicker of freedom objection holds that Frankfurt cases leave the agent an alternative after all: she could have shown the sign, or begun the deliberation, that would have triggered the intervener, and responsibility may require no more than this. We argue that the flicker cannot do the work the objection needs. Either the flicker is involuntary, in which case it is not something the agent could have done to avoid responsibility, or it is voluntary, in which case buffer cases remove it. In neither case does responsibility depend on it, and the conclusion of ACT-02, that responsibility does not require alternative possibilities, stands.
\end{abstract}

\section{The objection}

ACT-02-a argues that some agent is responsible for an action although she had no robust alternative to it (Theorem~3.1), where an alternative is robust if the agent could voluntarily have taken it and taking it would have been relevant to her responsibility (Definition~2.2). The flicker objection grants that Frankfurt's intervener removes every robust alternative but points to what remains \citep{Fischer1994}. Had Jones been about to decide otherwise, he would have shown the sign on which Black's intervention depends. Jones could therefore have shown it, and this flicker, a sliver of alternative possibility, is what responsibility requires. A version of the thought underlies van Inwagen's principles about particular events: what Jones is responsible for is the event that occurred, and something else could have occurred instead \citep{vanInwagen1978}.

\section{The dilemma for the flicker}

\begin{proposition}\label{p:dilemma}
Every flicker in a Frankfurt case is either involuntary or removed by a buffer construction. In neither case is it a robust alternative.
\end{proposition}

\begin{proof}
Suppose first that the flicker is involuntary: a blush, a neural pattern, an inclination that precedes any decision. Then it is not something the agent could have done, in the sense of Definition~2.2 of ACT-02-a, and so not a robust alternative. Suppose instead that it is voluntary: the agent could have begun to deliberate otherwise, for instance by attending to the reason against her action. Then the buffer construction of ACT-02-a, Section~4, applies: the intervener is triggered by the agent's failing to meet a condition necessary for deliberating otherwise, so that the voluntary alternative is closed before the agent could take it. What remains is again an involuntary difference, and the first case applies.
\end{proof}

The proposition shows that the flicker theorist must hold that responsibility can depend on an alternative that is not robust. This is not credible. An alternative bears on responsibility because, if responsibility requires the ability to avoid blame, it must be an ability one could exercise with that aim. An involuntary flicker is not such an ability. Jones could not have blushed in order to avoid responsibility. Had he blushed, he would not have been blamed for the decision Black then produced, but the blush would not have been his way of avoiding blame; it would have been an accident of his physiology.

The point can be put in terms of the actual sequence. By Principle~P1 of ACT-02-a, responsibility depends only on the actual sequence of deliberation, decision and action. The flicker is a feature of a non-actual sequence, and whether it is available depends on the design of the intervener's trigger, a factor idle in what Jones actually does. Two agents with the same actual sequence, one facing an intervener whose trigger is a blush and one facing an intervener whose trigger is a non-voluntary neural precursor without outward sign, differ in their flickers and not in anything that makes them responsible.

\begin{objection}[The flicker is the event]
Responsibility is for particular events, and a different event would have occurred had the intervener acted. So Jones could have brought it about that a different particular occurred.
\end{objection}

\begin{reply}
Had the intervener acted, the resulting decision would not have been Jones's doing, and he would not have brought it about in the sense relevant to responsibility. That a different particular would have occurred, by a process Jones did not control, gives him no voluntary route by which to avoid anything. Particular-event principles are therefore satisfied in Frankfurt cases only by alternatives that are not robust, and they inherit the dilemma of Proposition~\ref{p:dilemma}.
\end{reply}

\section{Conclusion}

The flicker of freedom is real in many Frankfurt cases, but it is not robust, and buffer cases show that it can always be made involuntary. Since responsibility could depend only on alternatives an agent could take in order to avoid blame, it does not depend on flickers. The objection thus leaves the conclusion of ACT-02 intact: responsibility does not require alternative possibilities \citep{Frankfurt1969, Pereboom2001}.

\begin{thebibliography}{9}
\bibitem[Fischer(1994)]{Fischer1994} Fischer, J.~M. (1994). \emph{The Metaphysics of Free Will: An Essay on Control}. Blackwell.
\bibitem[Frankfurt(1969)]{Frankfurt1969} Frankfurt, H.~G. (1969). Alternate possibilities and moral responsibility. \emph{The Journal of Philosophy}, 66(23), 829--839.
\bibitem[Pereboom(2001)]{Pereboom2001} Pereboom, D. (2001). \emph{Living Without Free Will}. Cambridge University Press.
\bibitem[van Inwagen(1978)]{vanInwagen1978} van Inwagen, P. (1978). Ability and responsibility. \emph{The Philosophical Review}, 87(2), 201--224.
\end{thebibliography}

\end{document}
